\documentclass[12pt]{article}
\usepackage[utf8]{inputenc}
\usepackage[T1]{fontenc}
\usepackage{hyperref}
\title{Lachiguiri Zapotec Language Resources for K-12 Curriculum Development}
\author{Compiled by Claude Sonnet 5.5 (claude-sonnet-5-5) for Andr\'as Kornai}
\date{}
\begin{document}
\maketitle

\section{Introduction}
Lachiguiri Zapotec (ISO 639-3 zpa; Glottocode lach1250) is a Zapotecan language of the Isthmus of Oaxaca, Mexico. It is also called Northwestern Tehuantepec Zapotec or Santiago Lachiguiri Zapotec. Do not confuse it with Lachix\'io Zapotec (zpl) or Lachirioag Zapotec (zty). The search found very little material. This report is short because of that.

\section{Primary Educational Resources}
I found no primers, textbooks, or school programmes for this language.

\section{Linguistic and Reference Materials}
\begin{itemize}
\item Wikipedia, ``Lachiguiri Zapotec'' (a stub article): \url{https://en.wikipedia.org/wiki/Lachiguiri_Zapotec}
\item Glottolog entry lach1250: \url{https://glottolog.org/resource/languoid/id/lach1250}
\end{itemize}
No dictionary or grammar turned up in the search.

\section{Religious and Community Resources}
\begin{itemize}
\item Global Recordings Network, Zapotec, Lachiguiri (Bible stories and evangelistic audio): \url{https://globalrecordings.net/en/language/zpa}
\end{itemize}

\section{Digital and Audio Resources}
Global Recordings offers downloadable MP3 and slideshow archives for this language, through the page listed above.

\section{Curriculum Development Context}
The language is reportedly used by all in its community, but no teaching materials were found. A curriculum would need new materials developed with the community. The Isthmus Zapotec literature may offer models, but I did not research it here.

\end{document}